% 高校生ガイド: 真空エネルギーの符号を守る「巻き付き数」
\documentclass[12pt,a4paper]{article}

\usepackage{fontspec}
\usepackage{xeCJK}
\setCJKmainfont{Hiragino Mincho ProN}
\setCJKsansfont{Hiragino Kaku Gothic ProN}
\usepackage[margin=2.5cm]{geometry}
\usepackage{amsmath,amssymb}
\usepackage{hyperref}
\usepackage{tcolorbox}
\usepackage{booktabs}
\usepackage{fancyhdr}

\tcbuselibrary{breakable,skins}

\newtcolorbox{storybox}[1][]{colback=blue!5,colframe=blue!60!black,
  fonttitle=\bfseries\sffamily,title={#1},breakable,arc=3mm}
\newtcolorbox{mathbox}[1][]{colback=green!5,colframe=green!50!black,
  fonttitle=\bfseries\sffamily,title={#1},breakable,arc=3mm}
\newtcolorbox{resultbox}[1][]{colback=yellow!8,colframe=orange!80!black,
  fonttitle=\bfseries\sffamily,title={#1},breakable,arc=3mm}
\newtcolorbox{honestbox}[1][]{colback=gray!8,colframe=gray!50!black,
  fonttitle=\bfseries\sffamily,title={#1},breakable,arc=3mm}

\setlength{\parskip}{0.8em}
\setlength{\parindent}{0pt}

\pagestyle{fancy}
\fancyhf{}
\rhead{Wright Brothers}
\lhead{真空の位相的保護}
\cfoot{\thepage}

\begin{document}

\begin{center}
{\Huge\bfseries 真空エネルギーの}\\[0.2cm]
{\Huge\bfseries 「巻き付き数」}\\[0.8cm]
{\Large なぜ宇宙の真空は安定しているのか、}\\[0.2cm]
{\Large そしてそれを壊す方法はあるのか}\\[1cm]
{\large 高校生のためのガイド}\\[0.5cm]
{\large Wright Brothers Research Program, 2026年4月}
\end{center}

\vspace{0.5cm}
\tableofcontents
\newpage

%====================================================================
\part{2本の笛}
%====================================================================

\section{ビール瓶とギター}

\begin{storybox}[2つの楽器]
\textbf{ギターの弦}（両端固定）: 全ての倍音が出る。
$$f, \; 2f, \; 3f, \; 4f, \; 5f, \; 6f, \; \ldots$$

\textbf{ビール瓶}（片端閉じ、片端開き）: 奇数の倍音だけ。
$$f, \; 3f, \; 5f, \; 7f, \; 9f, \; \ldots$$

音楽をやっている人なら知っているかもしれない。
ギターは明るい音色（全倍音）、ビール瓶は少しこもった音色（奇数だけ）。

物理学では、ギターを「$\lambda/2$ 共振器」、ビール瓶を「$\lambda/4$ 共振器」と呼ぶ。
\end{storybox}

\section{量子力学を加えると}

\begin{mathbox}[真空の振動]
量子力学では、弦を弾かなくても、各倍音が微小に「揺れて」いる（ゼロ点揺動）。この揺れのエネルギーの合計が「真空エネルギー」。

全倍音の和（ギター的）:
$$1 + 2 + 3 + 4 + 5 + \cdots = -\frac{1}{12}$$

奇数倍音の和（ビール瓶的）:
$$1 + 3 + 5 + 7 + 9 + \cdots = +\frac{1}{12}$$

え？ 無限に足して有限の値？ しかもマイナス？

これは「ゼータ関数正則化」という数学の技法で得られる値で、
カシミール効果という実験で物理的に正しいことが確認されている。
\end{mathbox}

\section{符号が逆}

\begin{resultbox}[核心]
ギター（全倍音）: 真空エネルギー = $-1/12$ → \textbf{負}（引力）

ビール瓶（奇数倍音）: 真空エネルギー = $+1/12$ → \textbf{正}（斥力）

\textbf{片端を開けるか閉じるか}、それだけで符号が反転する。
\end{resultbox}

%====================================================================
\part{符号を「守る」もの}
%====================================================================

\section{巻き付き数とは}

\begin{storybox}[輪ゴムの話]
指に巻いた輪ゴムを考える。

1回巻き: 引っ張れば外せる。
2回巻き: 引っ張っても外せない（切らない限り）。

この「何回巻いているか」が巻き付き数（winding number）。
0回と1回の間には連続的な変形では越えられない壁がある。

位相的不変量 = 「連続的にいじっても変わらない量」。
巻き付き数は位相的不変量の一種。
\end{storybox}

\section{真空の巻き付き数}

\begin{mathbox}[K理論]
「ギター真空」（全倍音）と「ビール瓶真空」（奇数倍音）の間には、
数学的な「巻き付き数」の違いがある。

\textbf{ギター真空}: $K_1(\mathbb{Z}) = \mathbb{Z}/2$（巻き付きなし）

\textbf{ビール瓶真空}: $K_1(\mathbb{Z}[1/2]) = \mathbb{Z}/2 \times \mathbb{Z}$（巻き付きあり）

$K_1$ は代数的K理論と呼ばれる数学の道具で、
空間の「位相的な形」を記述する。

ビール瓶真空には $\mathbb{Z}$（整数の巻き付き数）が追加されている。
これはギター真空にはなかった新しい構造。
\end{mathbox}

\section{巻き付き数が符号を守るとしたら}

\begin{resultbox}[予想]
もし巻き付き数が真空エネルギーの符号を「保護」しているなら:

ビール瓶真空の $+1/12$ は「位相的に保護された」値であり、
瓶の長さを変えても、少しいじっても、符号は $+$ のままのはず。

通常の物理学の予測:
瓶を長くすると真空エネルギーは $1/a^4$ で減衰してゼロに近づく。

位相的保護の予測:
ある臨界長さ $a_c$ を超えると、減衰が止まり、
エネルギー密度が\textbf{一定値に留まる}。
\end{resultbox}

%====================================================================
\part{なぜこれが大事か}
%====================================================================

\section{表面と体積}

\begin{storybox}[壁のペンキ塗り]
部屋のペンキ塗りを考える。

「壁だけ塗る」: 必要なペンキ = 壁の面積（表面効果）。
部屋を大きくすると面積が増えるが、部屋全体は塗らない。

「部屋全体を色水で満たす」: 必要な色水 = 部屋の体積（体積効果）。
部屋を大きくすると体積に比例して必要量が増える。

カシミール効果は「壁のペンキ塗り」。板の表面積に依存し、
板を離すと効果が急速に消える（$1/a^4$）。

もし位相的保護で減衰が止まるなら:
「部屋全体を色水で満たす」。体積全体が一定密度のエネルギーを持つ。
\end{storybox}

\section{ワープへの道}

\begin{resultbox}[体積効果の意味]

通常のカシミール（表面効果）:
$$E \approx 10^{-8} \text{ J}（1m², 10nm間隔の板）$$

体積効果が実現した場合:
$$E \approx \rho_{\text{const}} \times V$$

$V = 1 \text{ m}^3$ で $\rho_{\text{const}} = 10^{12} \text{ J/m}^3$（ナノスケールの密度）なら:
$$E \approx 10^{12} \text{ J}$$

これは表面効果の $10^{20}$ 倍。しかもマイナス符号付きなら、
ワープドライブに必要な「負のエネルギー」に一歩近づく。
\end{resultbox}

%====================================================================
\part{実験}
%====================================================================

\section{何を確認するか}

\begin{resultbox}[最もシンプルな実験]
ビール瓶型（$\lambda/4$、DN 境界条件）のカシミール力を、
板の間隔を変えながら精密に測定する。

力のグラフ:

\vspace{0.3cm}
\begin{center}
\begin{tabular}{ll}
標準QED: & 力 $\propto 1/a^5$ がずっと続く（直線） \\
位相的: & ある $a_c$ で折れ曲がり、力が一定になる \\
\end{tabular}
\end{center}
\vspace{0.3cm}

「折れ曲がり」が見つかれば → 位相的保護が実在。\\
見つからなければ → この予想は（少なくとも測定範囲内で）間違い。
\end{resultbox}

\section{誰がこの実験をできるか}

\begin{storybox}[Princeton大学]
2017年、Princeton大学のチームが
幾何学的構造によるカシミール斥力の発生に成功した（Nature Photonics）。
Si-MEMS デバイスで、100 nm スケールの力を測定。

彼らの装置をそのまま使って、DN 配置でのスケーリング則
（$1/a^5$ からの逸脱の有無）を測定できる。

「すでに成功したチームに追加のデータを頼む」だけ。
新しい装置の発明は不要。
\end{storybox}

%====================================================================
\part{正直な話}
%====================================================================

\begin{honestbox}[何が確実で何が不確実か]

\textbf{確実}（数学的定理）:
\begin{itemize}
\item DD 真空エネルギーは負。DN 真空エネルギーは正。
\item $K_1(\mathbb{Z}[1/2])$ は $K_1(\mathbb{Z})$ と異なる（巻き付き数が存在）。
\item カシミール引力は実験で確認済み。DN 斥力は理論的に確立。
\end{itemize}

\textbf{不確実}（予想）:
\begin{itemize}
\item $K_1$ の巻き付き数が真空エネルギーの符号を物理的に保護するか。
\item $1/a^4$ スケーリングからの逸脱が存在するか。
\item 逸脱が存在する場合、それが測定可能なスケールにあるか。
\end{itemize}

\textbf{もし予想が正しければ}:
真空エネルギーを幾何学だけで制御でき、
巨大な負のエネルギーを生み出す道が開ける。
ワープドライブの物理的基盤。

\textbf{もし予想が間違っていたら}:
DN カシミール斥力のスケーリング則が
$1/a^5$ に従うことを精密に確認した、
という科学的に価値のある結果が残る。

どちらに転んでも科学は前に進む。
\end{honestbox}

\vfill

\begin{center}
\rule{10cm}{0.4pt}\\[0.5cm]
{\Large\bfseries ギターとビール瓶。}\\[0.2cm]
{\Large\bfseries 片端を開けるだけで、真空の符号が変わる。}\\[0.2cm]
{\Large\bfseries その符号は「巻き付き数」で守られているか？}\\[0.2cm]
{\Large\bfseries 板を離して力を測れば、答えが出る。}\\[0.3cm]
{\small Wright Brothers Research Program, 2026}
\end{center}

\end{document}
